\documentclass[../../../include/open-logic-section]{subfiles}

\begin{document}

\olfileid{sth}{story}{urelements}
\olsection{Urelemen utawa Ora?}

Ing sawetara bab sabanjure, kita bakal nyoba njupuk aksioma
saka konsepsi kumulatif-iteratif tumrap himpunan. Nanging
sadurunge nerusake, kita kudu ngandharake luwih akeh babagan
\emph{urelemen}.

Gambaran ing \olref[approach]{sec} mung ngidini kita nggawe
himpunan anyar saka \emph{himpunan} lawas. Nanging kita bisa
uga ngidini \emph{barang dudu himpunan}---sapi, babi,
butiran wedhi, utawa liyane---dadi !!{element}s himpunan.
Ing kasus iku, kita miwiti saka sawetara unsur dhasar,
yaiku \emph{urelemen}, banjur ing saben tataran $S$
kita nggawe ``kabeh himpunan kang bisa'' kanthi anggota
urelemen utawa himpunan kang kawangun ing tataran sadurunge
$S$ (ing kombinasi apa wae). Gambaran kang kawangun bakal
luwih kaya iki:
\begin{center}
	\begin{tikzpicture}[scale=0.6]
	\tikzset{cut_here/.style={densely dotted}}
	\draw (-4,6) -- (-1, 0) -- (1,0) -- (4,6); 
	\draw[cut_here] (-5,8)--(-4,6);
	\draw[cut_here] (5,8)--(4,6);
	\draw(-1.5, 1)--(1.5, 1);
	\draw(-2, 2)--(2, 2);
	\draw(-2.5, 3)--(2.5,3);
	\draw(-3, 4)--(3, 4);
	\draw(-3.5, 5)--(3.5, 5);
	\draw(-4, 6)--(4, 6);
	\node[label] at (2, 0) {\small 0};
	\node[label] at (2.5, 1) {\small 1};
	\node[label] at (3, 2) {\small 2};
	\node[label] at (3.5, 3) {\small 3};
	\node[label] at (4, 4) {\small 4};
	\node[label] at (4.5, 5) {\small 5};
	\node[label] at (5, 6) {\small 6};
	\end{tikzpicture}
\end{center}
Mula saiki kita kudu mutusake: \emph{Apa urelemen kudu diidini?}

Saka sisi filsafat, lumrah yen urelemen dilebokake ing teori kita.
Alesan utamane supaya teori himpunan bisa \emph{ditrapake}.
Kanggo nggambarake iki, elinga saka \olref[sfr][siz][]{chap}
yen rong himpunan $A$ lan~$B$ diarani padha gedhene, yaiku
$\cardeq{A}{B}$, yen lan mung yen ana bijeksi ing antarane
loro-lorone. Yen sapi-sapi ing pangonan lan babi-babi ing
kandhang saben-saben dadi himpunan, kita bisa nerangake
pratelan ``cacah sapi padha karo cacah babi'' nganggo teori
himpunan. Nanging yen urelemen dilarang, nganti sapi lan
babi \emph{ora} dadi himpunan, panjelasan teoretis-himpunan
kasebut ora kasedhiya. Malah kita ora duwe cara langsung
kanggo ngetrapake teori himpunan marang apa wae saliyane
himpunan dhewe. (Kanggo alasan liyane kanggo nglebokake
urelemen, delengen \citealt[pp.~vi, 24,
50--1]{Potter2004}.)

Nanging saka sisi matematika, ngidini urelemen iku arang.
Salah siji alesane yaiku ngrumusake teori himpunan tanpa
urelemen iku \emph{rada sithik} luwih gampang. Nanging
sok-sok ana alasan kang luwih narik kawigaten kanggo
ora nglebokake urelemen ing teori himpunan:
\begin{quote}
	Miturut kapercayan yen teori himpunan minangka dhasar
	matematika, kita kudune bisa nyakup kabeh matematika
	kanthi mung ngrembug himpunan, mula variabel kita ora
	kudu nglimputi objek kaya sapi lan babi.
	%Nanging yen $C$ iku sapi, $\{C\}$ iku himpunan, nanging dudu objek matematika kang sah.
	\citep[p.~8]{Kunen1980}
\end{quote}
Dadi perhatian marang katerapan nyaranake \emph{nglebokake}
urelemen; dene tujuan dhasar kang reduktif (nyuda
matematika dadi teori himpunan murni) bisa nyaranake
\emph{ora nglebokake} urelemen. Rasa males kang entheng
uga nuding marang pilihan ora nglebokake urelemen.

Kita bakal milih cara kang paling gampang. Nanging pilihan
iki uga nduweni alasan pedagogis. Tujuane yaiku ngenalake
unsur-unsur teori himpunan kang dibutuhake kanggo wiwit
nyinaoni filsafat teori himpunan. Lan akeh-akehe pustaka
filsafat iku ngrembug teori himpunan kang dirumusake
\emph{tanpa} urelemen. Mula buku iki mbokmenawa luwih
migunani yen cukup cedhak karo pustaka kasebut.

\end{document}
